\section*{第 \ref{ch:b1:matrices} 章——矩阵}

\begin{solution}{exo:b1:matrices:1}
\[
AB = \begin{pmatrix} 2 & 1\\ 1 & 0\end{pmatrix},
\quad
BA = \begin{pmatrix} 0 & 1\\ 1 & 2\end{pmatrix},
\quad
A^2 - B^2 = \begin{pmatrix} 1 & 4\\ 0 & 1\end{pmatrix} - I
= \begin{pmatrix} 0 & 4\\ 0 & 0\end{pmatrix},
\]
\[
(A+B)(A-B) = A^2 - AB + BA - B^2
= \begin{pmatrix} 0 & 4\\ 0 & 0\end{pmatrix} +
\begin{pmatrix} -2 & 0\\ 0 & 2 \end{pmatrix}
= \begin{pmatrix} -2 & 4\\ 0 & 2\end{pmatrix}.
\]
二者之差为 $BA - AB \neq 0$：恒等式 $(a+b)(a-b) = a^2 -
b^2$ 要求乘法满足交换律，此处不成立。
\end{solution}

\begin{solution}{exo:b1:matrices:2}
$M$：$L_2 \leftarrow L_2 - 2L_1$ 使第二行化为零；$L_3
\leftarrow L_3 - L_1$ 得 $(0, -1, -2)$。有两个主元：
$\operatorname{rk} M = 2$。

$N$：$L_3 \leftarrow L_3 - L_1$ 得 $(0,1,1,1) = L_2$；再由 $L_3
\leftarrow L_3 - L_2 = 0$，有两个主元：$\operatorname{rk} N = 2$。
\end{solution}

\begin{solution}{exo:b1:matrices:3}
对 $(A \mid I_3)$ 作行化简：$L_2 \leftarrow L_2 - 2L_1$，$L_3
\leftarrow L_3 - L_1$：
\[
\begin{pmatrix}
1 & 0 & 1 & 1 & 0 & 0\\
0 & 1 & -1 & -2 & 1 & 0\\
0 & 1 & 0 & -1 & 0 & 1
\end{pmatrix}
\xrightarrow{L_3 \leftarrow L_3 - L_2}
\begin{pmatrix}
1 & 0 & 1 & 1 & 0 & 0\\
0 & 1 & -1 & -2 & 1 & 0\\
0 & 0 & 1 & 1 & -1 & 1
\end{pmatrix},
\]
再由 $L_1 \leftarrow L_1 - L_3$，$L_2 \leftarrow L_2 + L_3$ 得：
\[
A^{-1} = \begin{pmatrix}
0 & 1 & -1\\
-1 & 0 & 1\\
1 & -1 & 1
\end{pmatrix}.
\]
\emph{验算：}$A$ 的第一行乘 $A^{-1}$ 的第一列：$1
\cdot 0 + 0\cdot(-1) + 1\cdot 1 = 1$；乘第二列：$1 - 0 -
1 = 0$；乘第三列：$-1 + 0 + 1 = 0$。
\end{solution}

\begin{solution}{exo:b1:matrices:4}
$u(1) = 1$，$u(X) = X + 1$，$u(X^2) = X^2 + 2X + 1$：以各像的
在 $(1, X, X^2)$ 下的\omterm{prop:b1:vspaces:coordinates}{坐标}为列，得
\[
M = \begin{pmatrix}
1 & 1 & 1\\
0 & 1 & 2\\
0 & 0 & 1
\end{pmatrix}.
\]
$u$ 可逆，因其有显然的逆 $P \mapsto P(X
- 1)$（即代换的复合）。其矩阵同样由
$u^{-1}(X^k) = (X-1)^k$ 求得：
\[
M^{-1} = \begin{pmatrix}
1 & -1 & 1\\
0 & 1 & -2\\
0 & 0 & 1
\end{pmatrix}.
\]
\end{solution}

\begin{solution}{exo:b1:matrices:5}
$N = \begin{pmatrix} 0 & 1\\ 0 & 0\end{pmatrix}$，$N^2 = 0$。由于
$2I$ 与 $N$ 可交换，二项式展开只保留前两项，
于是：
\[
A^k = (2I + N)^k = 2^k I + k\,2^{k-1} N
= \begin{pmatrix} 2^k & k\,2^{k-1}\\ 0 & 2^k\end{pmatrix}.
\]
（验算 $k = 2$：$A^2 = \begin{pmatrix}4 & 4\\ 0 & 4\end{pmatrix}$，
与直接乘法的结果一致。）
\end{solution}

\begin{solution}{exo:b1:matrices:6}
迹：$\operatorname{tr}(AB - BA) = \operatorname{tr}(AB) -
\operatorname{tr}(BA) = 0$（见\cref{def:b1:matrices:transpose}），
而 $\operatorname{tr}(I_n) = n \neq 0$（在 $\R$ 或 $\C$ 中），因此
无解。（在无限维空间中，恒等式\emph{确实可以}
实现——求导与乘以 $x$ 满足它
——正是因为那里不存在迹。）
\end{solution}

\begin{solution}{exo:b1:matrices:7}
裂项相乘，$A$ 的各次幂两两可交换：
\[
(I - A)(I + A + \dots + A^{m-1}) = I - A^m = I ,
\]
由\cref{prop:b1:matrices:ring}可知单侧逆即是逆。至于应用：
所给矩阵为 $I + N$，其中
\[
N = \begin{pmatrix} 0 & 2 & 3\\ 0 & 0 & 2\\ 0&0&0 \end{pmatrix},
\quad
N^2 = \begin{pmatrix} 0&0&4\\ 0&0&0\\ 0&0&0\end{pmatrix},
\quad N^3 = 0 ,
\]
于是在公式中将 $A$ 换为 $-N$：
\[
(I + N)^{-1} = I - N + N^2 =
\begin{pmatrix}
1 & -2 & 1\\
0 & 1 & -2\\
0 & 0 & 1
\end{pmatrix}.
\]
\end{solution}

\begin{solution}{exo:b1:matrices:8}
$A^2 = A$：自\omterm{def:b1:structures:morphism}{同态} $a$ 是\omterm{def:b1:linmaps:projection}{投影}
（由\cref{thm:b1:linmaps:projchar}），$E = \operatorname{im} a \oplus
\ker a$，且 $\dim\operatorname{im} a = r = \operatorname{rk} A$。在
与该分解相适应的基（取像中的 $r$ 个向量，
再取核的一组基）下，$a$ 的矩阵为
$\begin{pmatrix} I_r & 0\\ 0 & 0\end{pmatrix}$，迹为 $r$。
迹在\omterm{def:b1:matrices:changeofbasis}{基变换}下保持不变：
$\operatorname{tr}(P^{-1}MP) = \operatorname{tr}(MPP^{-1}) =
\operatorname{tr} M$（由循环恒等式）。因此
$\operatorname{tr} A = r = \operatorname{rk} A$.
\end{solution}

\begin{solution}{exo:b1:matrices:9}
$J^2 = nJ$（$J^2$ 的每个元素是 $n$ 个 1 之和）。设 $M^{-1} =
\alpha I + \beta J$:
\[
(aI + bJ)(\alpha I + \beta J)
= a\alpha\, I + (a\beta + b\alpha + nb\beta)\, J .
\]
该乘积等于 $I$ 当且仅当 $a\alpha = 1$ 且 $a\beta + b\alpha + nb\beta =
0$，即 $\alpha = \frac1a$ 且 $\beta(a + nb) = -\frac ba$。若 $a
\neq 0$ 且 $a + nb \neq 0$，则
\[
M^{-1} = \frac 1a I - \frac{b}{a(a + nb)}\, J .
\]
反之，若 $a = 0$，则 $M = bJ$ 的秩 $\leq 1 < n$（当 $n \geq
2$）：不可逆（$n = 1$ 为数量情形）。若 $a + nb = 0$：
向量 $v = (1, \dots, 1)^{\mathsf T}$ 满足 $Mv = (a +
nb)v = 0$，且 $v \neq 0$，故不可逆。因此 $M \in GL_n \iff a
\neq 0$ 且 $a + nb \neq 0$。
\end{solution}

\begin{solution}{exo:b1:matrices:10}
\emph{和：}$\operatorname{im}(A + B) \subseteq \operatorname{im} A
+ \operatorname{im} B$（因 $(A+B)x = Ax + Bx$），由格拉斯曼
公式可知，和的维数不超过各维数之和。

\emph{西尔维斯特：}设 $a$ 为 $A$ 限制于 $V =
\operatorname{im} B$（维数为 $\operatorname{rk} B$）。其像为
$\operatorname{im}(AB)$（$a(Bx) = ABx$），并在 $V$ 上应用秩--零化度定理：
\[
\operatorname{rk} B = \dim\ker(a_{|V}) + \operatorname{rk}(AB) .
\]
现在 $\ker(a_{|V}) \subseteq \ker A$，其维数为 $n -
\operatorname{rk} A$，于是
\[
\operatorname{rk}(AB) \geq \operatorname{rk} B - (n -
\operatorname{rk} A) = \operatorname{rk} A + \operatorname{rk} B -
n . \qedhere
\]
\end{solution}

\begin{solution}{exo:b1:matrices:11}
\begin{enumerate}
  \item 逐元素比较，$(AD)_{ij} = a_{ij}\,d_j$，$(DA)_{ij} =
        d_i\,a_{ij}$。即 $AD = DA$ 当且仅当 $a_{ij}(d_j - d_i) = 0$
        对所有 $i, j$ 成立；当 $i \neq j$ 时，因子 $d_j - d_i$
        非零，从而 $a_{ij} = 0$，$A$ 为对角矩阵。
        反之，对角矩阵两两可交换。
  \item 若 $A$ 与所有矩阵可交换，则它与
        $\operatorname{diag}(1, 2, \dots, n)$ 可交换，故 $A =
        \operatorname{diag}(\lambda_1, \dots, \lambda_n)$，由（1）。
        则 $A E_{ij} = \lambda_i E_{ij}$（仅第 $i$ 行
        的 $E_{ij}$ 保留），而 $E_{ij} A = \lambda_j E_{ij}$：
        由与 $E_{ij}$ 可交换，得 $\lambda_i = \lambda_j$。
        于是 $A = \lambda I_n$。数量矩阵与一切矩阵
        都可交换，故 $\mathcal{M}_n(K)$ 的中心为
        $K\,I_n$.
\end{enumerate}
\end{solution}

\begin{solution}{exo:b1:matrices:12}
\begin{enumerate}
  \item 若 $\operatorname{rk} A = 1$，则 $A$ 的像为一条直线
        $\operatorname{Vect}(C)$，$C \neq 0$，故对第 $j$ 列，
        $A$ 的该列为 $\ell_j\,C$，其中 $\ell_j$ 为标量（不全为
        零），即 $A = C L$，其中 $L = (\ell_1, \dots,
        \ell_n) \neq 0$。反之，若 $A = CL \neq 0$，则所有
        列都是 $C$ 的倍数：秩为 $1$。
  \item $A^2 = C\,(L C)\,L$，而 $LC$ 为数量 $\sum_i
        \ell_i c_i = \operatorname{tr}(CL) = \operatorname{tr}
        A$。于是 $A^2 = (\operatorname{tr} A)\,A$，由
        归纳法得 $A^m = (\operatorname{tr} A)^{m-1} A$。若
        $\operatorname{tr} A \neq 0$，则任何幂都不为零；若
        $\operatorname{tr} A = 0$，则 $A^2 = 0$。秩为一的
        矩阵幂零当且仅当其迹为零。
  \item 记 $t = \operatorname{tr} A \neq -1$：
        \[
        (I_n + A)\Bigl(I_n - \frac{A}{1 + t}\Bigr)
        = I_n + A - \frac{A + A^2}{1 + t}
        = I_n + A - \frac{(1 + t)A}{1 + t} = I_n ,
        \]
        利用 $A^2 = tA$。若 $t = -1$，则 $(I_n + A)A = A + A^2 = A
        - A = 0$，且 $A \neq 0$，故 $I_n + A$ 零化
        $A$ 的每个（非零）列：不是\omterm{def:b1:logic:inj}{单射}，从而不可逆。
\end{enumerate}
\end{solution}

\begin{solution}{pb:b1:matrices:1}
\textbf{1.} 对 $X^n$ 作欧几里得除法，除式为 $2$ 次\omterm{def:b1:poly:def}{首一多项式}
$D$（见\cref{thm:b1:poly:division}）：商式与余式存在
且唯一，余式次数 $\leq 1$：$X^n =
Q_n D + a_n X + b_n$。当 $n = 0$ 时，$Q_0 = 0$，$(a_0, b_0) = (0,
1)$；当 $n = 1$ 时，$(a_1, b_1) = (1, 0)$。

\textbf{2.} 乘以 $X$，并利用 $X^2 = D + sX - p$ 降次：
\[
X^{n+1} = X Q_n D + a_n X^2 + b_n X
= (X Q_n + a_n)\,D + (s\,a_n + b_n)\,X - p\,a_n .
\]
最后一式具有余式的形式（次数 $\leq 1$），由
唯一性，$a_{n+1} = s a_n + b_n$，$b_{n+1} = -p a_n$。
将 $b_{n+1} = -pa_n$ 代入 $a_{n+2} = s a_{n+1} +
b_{n+1}$，得 $a_{n+2} = s\,a_{n+1} - p\,a_n$。

\textbf{3.} 将 $X^n = Q_n D + a_n X + b_n$ 在两根处取值：
$\lambda^n = a_n\lambda + b_n$ 与 $\mu^n = a_n\mu + b_n$。
两式相减，并除以 $\lambda - \mu \neq 0$：
\[
a_n = \frac{\lambda^n - \mu^n}{\lambda - \mu},
\qquad
b_n = \lambda^n - a_n\lambda
= \frac{\lambda\mu^n - \mu\lambda^n}{\lambda - \mu} .
\]

\textbf{4.} 在二重根处：$\lambda^n = a_n\lambda + b_n$。
对恒等式求导，得 $nX^{n-1} = Q_n'\,(X - \lambda)^2
+ 2Q_n\,(X - \lambda) + a_n$，并在 $\lambda$ 处取值：$a_n
= n\lambda^{n-1}$；于是 $b_n = \lambda^n - n\lambda^{n} = (1 -
n)\lambda^{n}$.

\textbf{5.} 对 $P = \sum_i p_i X^i$ 与 $Q = \sum_j q_j X^j$，
\[
P(M)\,Q(M) = \sum_{i,j} p_i q_j M^{i+j} = (PQ)(M),
\]
因为单一矩阵 $M$ 的各次幂彼此可交换
（和式由线性显然成立）。若 $D(M) = 0$，将 $M$
代入 $X^n = Q_n D + a_n X + b_n$，得 $M^n = Q_n(M)\,D(M) + a_n
M + b_n I = a_n M + b_n I$.

\textbf{6.} 直接计算乘积：
\[
A^2 = \begin{pmatrix}
a^2 + bc & b(a + d)\\
c(a + d) & d^2 + bc
\end{pmatrix},
\qquad
s A = \begin{pmatrix}
a(a+d) & b(a+d)\\
c(a+d) & d(a+d)
\end{pmatrix},
\]
于是 $A^2 - sA$ 的非对角元全为零，对角元为
$a^2 + bc - a^2 - ad = bc - ad = -p$: $A^2 - sA + pI_2 = 0$.

\textbf{7.} 设 $A' = \begin{pmatrix} a' & b'\\ c' &
d'\end{pmatrix}$，展开得 $p(AA') = (aa' + bc')(cb' + dd') -
(ab' + bd')(ca' + dc')$：项 $aa'cb'$ 与 $ab'ca'$ 相消，
项 $bc'dd'$ 与 $bd'dc'$ 相消，余下为
\[
aa'dd' - bca'd' + bcb'c' - adb'c'
= (ad - bc)(a'd' - b'c') = p(A)\,p(A').
\]
若 $p \neq 0$，由凯莱--哈密顿定理得 $A\,\bigl(\tfrac1p(sI_2 -
A)\bigr) = \tfrac1p(sA - A^2) = I_2$，从而得逆矩阵（
由\cref{prop:b1:matrices:ring}可知它是双侧的）。若 $p = 0$ 且
$A$ 可逆，则由乘法性得 $1 = p(I_2) =
p(A)\,p(A^{-1}) = 0$：矛盾。故 $A \in GL_2 \iff p \neq 0$。

\textbf{8.} 取 $s = 3$，$p = 2$，$D = X^2 - 3X + 2 = (X - 1)(X -
2)$：$\lambda = 2$，$\mu = 1$，于是 $a_n = 2^n - 1$，且 $b_n = 2 -
2^n$（第 3 题）。由此
\[
A^n = (2^n - 1)A + (2 - 2^n)I
= \begin{pmatrix} 1 & 2^n - 1\\ 0 & 2^n \end{pmatrix}.
\]
验算：$A^2 = \begin{pmatrix} 1 & 3\\ 0 & 4\end{pmatrix}$ 既可由公式给出，
也可直接平方算出。

\textbf{9.} 取 $s = 4$，$p = 3\cdot1 - 1\cdot(-1) = 4$：$D = X^2 -
4X + 4 = (X - 2)^2$，为二重根 $\lambda = 2$。由第 4 题，$a_n
= n\,2^{n-1}$，$b_n = (1 - n)2^n$，于是
\[
A^n = n\,2^{n-1}A + (1 - n)2^n I
= 2^{n-1}\begin{pmatrix} n + 2 & n\\ -n & 2 - n
\end{pmatrix}.
\]
当 $n = 2$ 时：$2\begin{pmatrix} 4 & 2\\ -2 & 0\end{pmatrix} =
\begin{pmatrix} 8 & 4\\ -4 & 0 \end{pmatrix}$，即 $A^2$
直接计算所得。

\textbf{10.} 归纳法：$F^1 = \begin{pmatrix} F_2 & F_1\\ F_1 &
F_0\end{pmatrix}$，且
\[
F^{n+1} = F^n F =
\begin{pmatrix} F_{n+1} + F_n & F_{n+1}\\
F_n + F_{n-1} & F_n \end{pmatrix}
= \begin{pmatrix} F_{n+2} & F_{n+1}\\ F_{n+1} & F_n
\end{pmatrix}.
\]
此处 $s = 1$，$p = -1$，$D = X^2 - X - 1$，其根为 $\varphi,
\psi$（$\varphi - \psi = \sqrt5$，$\varphi\psi = -1$）。
数列 $(F_n)$ 满足 $F_0 = 0 = a_0$，$F_1 = 1 = a_1$，且服从
与 $(a_n)$ 相同的递推式：$F_n = a_n = (\varphi^n -
\psi^n)/\sqrt5$，即比内公式。卡西尼恒等式：应用第 7 题的
乘法性于 $F^n$，有
\[
F_{n+1}F_{n-1} - F_n^2 = p(F^n) = p(F)^n = (-1)^n .
\]

\textbf{11.} 该条件是线性的，且包含零数列，
故为\omterm{def:b1:vspaces:subspace}{子空间}。由归纳法，$u_0, u_1$ 线性地决定 $u$，
且每一对初值都恰好对应
唯一一个解：仿\cref{exo:b1:findim:10}，$E_D$
可由 $(u_0, u_1) \in K^2$ \omterm{def:b1:logic:inj}{双射}且线性地参数化：
$\dim E_D = 2$.

\textbf{12.} $(a_n)$ 满足（第 2 题）的递推式，其中 $a_0
= 0$，$a_1 = 1$。$(b_n)$ 亦然：$b_{n+2} = -p\,a_{n+1} =
-p(s a_n + b_n) = s\,b_{n+1} - p\,b_n$（两次利用 $b_{n+1} = -pa_n$
），且 $b_0 = 1$，$b_1 = 0$。组合 $v_n =
u_1 a_n + u_0 b_n$ 亦为解，且 $v_0 = u_0$，$v_1 =
u_1$；初值相同的两个解必相等，
（由归纳法），故 $u_n = u_1 a_n + u_0 b_n$ 对一切 $n$ 成立。

\textbf{13.} $(\lambda^n)$ 为解当且仅当 $\lambda^{n+2} =
s\lambda^{n+1} - p\lambda^n$ 对所有 $n$ 成立，即 $D(\lambda) = 0$
（两边除以 $\lambda^n \neq 0$；注意 $\lambda, \mu \neq
0$，因为 $p = \lambda\mu \neq 0$）。证明
$\bigl((\lambda^n), (\mu^n)\bigr)$ \omterm{def:b1:vspaces:free}{线性无关}：在 $n = 0, 1$ 处的关系式
给出 $c + c' = 0$，$c\lambda + c'\mu = 0$，于是 $c(\lambda - \mu)
= 0$，得 $c = c' = 0$。在 $2$ 维空间中，两个\omterm{def:b1:vspaces:free}{线性无关}的向量即为基。
重根情形：$\bigl((n\lambda^n)\bigr)$ 为解，
因 $s = 2\lambda$，$p = \lambda^2$：
\[
s(n+1)\lambda^{n+1} - p\,n\lambda^n
= \lambda^{n+2}\bigl(2(n+1) - n\bigr) = (n+2)\lambda^{n+2} ;
\]
\omterm{def:b1:vspaces:free}{线性无关}性：在 $n = 0, 1$ 处得 $c = 0$，再由 $c'\lambda = 0$
且 $\lambda \neq 0$，于是两系数皆为零。

\textbf{14.} $D = X^2 - X - 6 = (X - 3)(X + 2)$。通解
为 $u_n = A\,3^n + B(-2)^n$；由初始条件得
$A + B = 1$，$3A - 2B = 8$，故 $A = 2$，$B = -1$：
\[
u_n = 2\cdot 3^n - (-2)^n .
\]
验算：$u_2 = 18 - 4 = 14 = u_1 + 6u_0$；$u_3 = 54 + 8 = 62 =
u_2 + 6u_1 = 14 + 48$。

\textbf{15.} $C\begin{pmatrix} u_n\\ u_{n+1}\end{pmatrix} =
\begin{pmatrix} u_{n+1}\\ -p\,u_n + s\,u_{n+1}\end{pmatrix} =
\begin{pmatrix} u_{n+1}\\ u_{n+2}\end{pmatrix}$，由归纳法
给出含 $C^n$ 的公式。又 $\operatorname{tr} C = 0
+ s = s$，且 $p(C) = 0\cdot s - 1\cdot(-p) = p$：友矩阵
恰以 $D$ 为其凯莱--哈密顿\omterm{def:b1:poly:def}{多项式}。

\textbf{16.} 设 $u$ 为满足 $u_{n+3} = \alpha u_{n+2} +
\beta u_{n+1} + \gamma u_n$ 的任一解，则状态向量 $v_n = (u_n,
u_{n+1}, u_{n+2})^{\mathsf T}$ 满足 $C_3 v_n = v_{n+1}$
（前两行整体平移，末行由
递推关系给出）。于是
\[
D_3(C_3)\,v_0 = v_3 - \alpha v_2 - \beta v_1 - \gamma v_0 ,
\]
其三个分量为 $u_{k+3} - \alpha u_{k+2} - \beta
u_{k+1} - \gamma u_k = 0$（$k = 0, 1, 2$）。
由于初始状态 $v_0 = (u_0, u_1, u_2)^{\mathsf T}$ 取遍\emph{全部}
$K^3$ 中的向量（初值可任意取定），矩阵 $D_3(C_3)$ 零化
每个向量，于是 $D_3(C_3) = 0$。

\textbf{17.} 写 $X^n = Q\,D_3 + R_n$，其中 $\deg R_n \leq 2$
并在各根处求值：$\lambda_i^n = R_n(\lambda_i)$。于是
$R_n$ 是次数 $\leq 2$ 的\omterm{def:b1:poly:def}{多项式}，它插值三个
值 $\lambda_i^n$（位于三个互异节点 $\lambda_i$ 处）。由
\cref{thm:b1:poly:lagrange}中的唯一性，$R_n = \sum_i
\lambda_i^n L_i$，其中 $(L_i)$ 为节点的拉格朗日基。
代入 $C_3$（第 5 题与第 16 题）：
\[
C_3^{\,n} = R_n(C_3) = \sum_{i=1}^{3} \lambda_i^n\,L_i(C_3),
\]
其中三个矩阵 $L_i(C_3)$ 与 $n$ 无关，故
$C_3^{\,n}$ 的每一个元素，都是 $\lambda_1^n,
\lambda_2^n, \lambda_3^n$ 的固定线性组合。

\textbf{18.} $D_3 = X^3 - 2X^2 - X + 2 = (X-1)(X+1)(X-2)$。
通解为 $u_n = A + B(-1)^n + C\,2^n$。初始
条件：$A + B + C = 0$，$A - B + 2C = 1$，$A + B + 4C = 1$。
将第三式减去第一式：$3C = 1$，$C = \frac13$；
于是 $A + B = -\frac13$，$A - B = \frac13$，得 $A = 0$，$B =
-\frac13$。由此
\[
u_n = \frac{2^n - (-1)^n}{3}
\]
（雅各布斯塔尔数）。验算：$u_3 = \frac{8 + 1}{3} = 3 =
2u_2 + u_1 - 2u_0 = 2 + 1 - 0$。

\textbf{19.} \omterm{def:b1:poly:def}{多项式} $X^n$ 在
$\lambda$ 处的泰勒展开：
\[
X^n = \sum_{k=0}^{n} \binom nk \lambda^{n-k}(X - \lambda)^k ,
\]
而 $k \geq 3$ 的各项均可被 $(X -
\lambda)^3$ \omterm{def:b1:arith:divides}{整除}，余式为
\[
R_n = \lambda^n + n\lambda^{n-1}(X - \lambda) + \binom
n2\lambda^{n-2}(X - \lambda)^2 .
\]
设 $M = \lambda I + N$，且 $N^3 = 0$，则 $(M - \lambda I)^3 = N^3
= 0$，由第 5 题可得
\[
M^n = \lambda^n I + n\lambda^{n-1} N + \binom n2
\lambda^{n-2} N^2 ,
\]
这恰是 $(\lambda I + N)^n$ 的二项展开
在 $N^2$ 处截断——两种方法结果一致。

\textbf{20.} 解空间的维数为 $3$（参数化方式与
第 11 题相同，由 $(u_0, u_1, u_2)$ 参数化），且
每个 $(\lambda_i^n)$ 都是解。\omterm{def:b1:vspaces:free}{线性无关}性：设
$c_1\lambda_1^n + c_2\lambda_2^n + c_3\lambda_3^n = 0$，对 $n =
0, 1, 2$ 成立。固定 $i$，记 $L_i = \sum_{k \leq 2} p_k X^k$ 为
关于节点的拉格朗日\omterm{def:b1:poly:def}{多项式}，满足 $L_i(\lambda_j) =
\delta_{ij}$。于是
\[
0 = \sum_{k=0}^{2} p_k\Bigl(\sum_j c_j\lambda_j^k\Bigr)
= \sum_j c_j\,L_i(\lambda_j) = c_i .
\]
故所有 $c_i = 0$，三个\omterm{def:b1:vspaces:free}{线性无关}的解在 $3$ 维空间中
构成一组基；通解为 $c_1\lambda_1^n +
c_2\lambda_2^n + c_3\lambda_3^n$。

\textbf{21.} 由 $F_k = F_{k+2} - F_{k+1}$，求和可裂项相消：
\[
\sum_{k=1}^{n} F_k = \sum_{k=1}^{n}\bigl(F_{k+2} - F_{k+1}\bigr)
= F_{n+2} - F_2 = F_{n+2} - 1 .
\]

\textbf{22.} 取 $(1,2)$ 元素，对 $F^{m+n} = F^m F^n$ 两端比较：
左边为 $F_{m+n}$；右边为第 $1$ 行（$F^m$ 的）
与第 $2$ 列（$F^n$ 的）的乘积，即 $F_{m+1}F_n + F_m F_{n-1}$。
令 $m = n$：
\[
F_{2n} = F_{n+1}F_n + F_nF_{n-1} = F_n\,(F_{n+1} + F_{n-1}).
\]

\textbf{23.} 由比内公式，$F_n - \dfrac{\varphi^n}{\sqrt5} =
-\dfrac{\psi^n}{\sqrt5}$，且 $\abs\psi = \frac{\sqrt5 - 1}2 <
1$，于是
\[
\Bigl|F_n - \frac{\varphi^n}{\sqrt5}\Bigr|
\leq \frac{1}{\sqrt5} < \frac12
\qquad (n \geq 0):
\]
$F_n$ 即为最接近 $\varphi^n/\sqrt5$ 的整数。

\textbf{24.} $t_n = F_{n+1} + F_{n-1}$ 是平移\omterm{pb:b1:matrices:1}{斐波那契数}列的线性组合，
故满足相同的递推式：
$t_{n+2} = t_{n+1} + t_n$；且 $t_1 = F_2 + F_0 =
1$，$t_2 = F_3 + F_1 = 3$。这些数即为卢卡斯数 $L_n$。
数列 $\varphi^n + \psi^n$ 也是该递推式的解，
且前两项相同（$\varphi + \psi = 1$，$\varphi^2 + \psi^2 = (
\varphi + \psi)^2 - 2\varphi\psi = 3$），故 $L_n = \varphi^n +
\psi^n$。最后
\[
F_n L_n = \frac{(\varphi^n - \psi^n)(\varphi^n +
\psi^n)}{\sqrt5} = \frac{\varphi^{2n} - \psi^{2n}}{\sqrt5} =
F_{2n},
\]
此即第 22 题的结论。

\textbf{25.} (i) 五个矩阵 $I, A, A^2, A^3, A^4$ 位于
$4$ 维空间 $\mathcal{M}_2(K)$ 中，故必有\emph{某个}非零
次数 $\leq 4$ 的\omterm{def:b1:poly:def}{多项式}零化 $A$；第二部分将其精确化为
一个显式的二次关系 $A^2 = sA - pI$，使所有幂次
都落在平面 $\operatorname{Vect}(I, A)$ 内。(ii)
用带余除法将 $X^n$ 对该二次式取余，
余式的两个系数满足二项递推关系，
$a_{n+2} = s\,a_{n+1} - p\,a_n$：幂运算已化为
迭代。(iii) 第 6 题是 $n = 2$ 时的
\emph{凯莱–哈密顿定理}的特例，在任意维数下均成立，
并已在本科二年级分册中证明。(iv) 伴侣矩阵使这一循环
闭合：每个\omterm{pb:b1:matrices:1}{线性递推}式\emph{即是}一个矩阵幂，其中
同一\omterm{def:b1:poly:def}{多项式} $D$ 以迹与行列式数据的形式出现，
从而余式演算可以
一步完成求解递推式与计算幂。
\end{solution}
